\subsection{跟算例：只约束到二阶时，最大熵分布为什么是高斯}
为把前面的抽象乘子变成可计算的例子，设已知数密度$n$、平均速度$\vel$和正定协方差
$\Theta_{ij}=P_{ij}/(mn)$。最大化$-\int f\ln(f/f_*)\dd^3v$，同时固定
$\int f$、$\int v_if$和$\int v_iv_jf$。对$f$变分后，驻值条件具有形式
\begin{equation}
 -\ln(f/f_*)-1+\lambda_0+\lambda_iv_i+\Lambda_{ij}v_iv_j=0.
\end{equation}
因此$f$是速度二次式的指数。为了使$\int f$有限，二次型$\Lambda$必须负定。把指数完成平方、用给定均值与协方差确定系数：
\begin{equation}
 f(\vct v)={n\over(2\pi)^{3/2}\sqrt{\det\Theta}}\,
 \exp\left[-{1\over2}(\vct v-\vel)^T\Theta^{-1}(\vct v-\vel)\right].
\end{equation}
若$\Theta=(T/m)\mathsf I$，便退回各向同性Maxwell分布。若$\Theta$有不同特征值，则沿不同主轴具有不同温度宽度。

该高斯关于$\vct c=\vct v-\vel$为偶函数，所以所有三阶中心矩
$Q_{ijk}=m\int c_ic_jc_kf\dd^3v$为零。这并不证明真实无碰撞重联的热流为零；它说明\emph{仅给二阶矩并选这一最大熵闭合}时，所重建分布没有偏斜和多束信息。若想保留非零三阶矩，就必须再加约束并处理更高阶可积性问题。

另一个实用结论是：若数值压力矩阵出现负特征值，$\Theta^{-1}$仍可能在代数上存在，但这份高斯已经不能作为物理概率密度。矩的“可实现性”就是要求存在非负、可积的分布产生这些矩，而不是仅要求程序能够对矩阵求逆。
